\begin{lemma}\label{lemgenn1}
Для последовательности из определения~\ref{defn1pil}, при $r=k+p-1$,
формальная производящая функция равна
\[
 \sum_{n\ge0}F_n^p(t,k)x^n
 =\frac{x^{r-1}}{1-tx-\sum_{j=p+1}^{r}x^j}.
\]
В частности, при $p=1$
\begin{equation}
 g_{F_n^1(t,k)}(x)
 =\sum_{n\ge0}F_n^1(t,k)x^n
 =\frac{x^{k-1}}{1-tx-x^2-\cdots-x^k}.
\end{equation}
\end{lemma}
\begin{proof}
Положим $G(x)=\sum_{n\ge0}F_n^p(t,k)x^n$ и
$D(x)=1-tx-\sum_{j=p+1}^{r}x^j$.
Соглашение о нулевых членах с отрицательными индексами даёт
\[
 [x^n](D(x)G(x))
 =F_n^p(t,k)-tF_{n-1}^p(t,k)
       -\sum_{j=p+1}^{r}F_{n-j}^p(t,k).
\]
При $n<r-1$ этот коэффициент равен $0$ по начальным условиям;
при $n=r-1$ он равен $1$; при $n\ge r$ он равен $0$
по рекурсии. Поэтому $D(x)G(x)=x^{r-1}$.
Свободный член $D$ равен $1$, следовательно, $D$ обратим
в кольце формальных степенных рядов. Это доказывает общую формулу.
При $p=1$ имеем $r=k$, что даёт указанное частное выражение.
\end{proof}
